\documentclass[12pt]{ximera}
\title{Homework 2: Sets, Subsets, and Set Operations}
\begin{document}
\begin{abstract}
Section 7.1: listing the elements of sets built from complements, differences, and
intersections; translating set identities into English; and counting the subsets of
a finite set.
\end{abstract}
\maketitle

\section*{Section 7.1: Sets, Subsets, and Set Operations}

\begin{problem}
Consider three sets
\[
A=\set{2,3,4,5}, \quad B=\set{3,5,7,9}, \quad C=\set{2,4,5,7,9}
\]
inside the universal set $\mathcal{U}=\set{1,2,3,4,5,6,7,8,9}$. Write out each of the
following sets, entering the elements in increasing order.

\begin{prompt}
\textbf{(a)} $\olsi{B} = \set{\answer{1},\ \answer{2},\ \answer{4},\ \answer{6},\ \answer{8}}$

\begin{feedback}
Removing $B=\set{3,5,7,9}$ from $\mathcal{U}=\set{1,\ldots,9}$ leaves
$\olsi{B}=\set{1,2,4,6,8}$.
\end{feedback}
\end{prompt}

\begin{prompt}
\textbf{(b)} $C - A = \set{\answer{7},\ \answer{9}}$

\begin{feedback}
$C=\set{2,4,5,7,9}$; removing the elements it shares with $A=\set{2,3,4,5}$ leaves
$C-A=\set{7,9}$.
\end{feedback}
\end{prompt}

\begin{prompt}
\textbf{(c)} $\olsi{B}\cap\left(C - A\right)$ has $\answer{0}$ elements.

\begin{multipleChoice}
\choice[correct]{It is the empty set $\varnothing$}
\choice{$\set{7}$}
\choice{$\set{7,9}$}
\choice{$\set{1,2,4,6,8}$}
\end{multipleChoice}

\begin{hint}
Compare the two sets you found in parts (a) and (b) element by element.
\end{hint}

\begin{feedback}
$\olsi{B}=\set{1,2,4,6,8}$ and $C-A=\set{7,9}$ have nothing in common, so their
intersection is the empty set: $\olsi{B}\cap(C-A)=\varnothing$. That makes sense:
$7$ and $9$ are both in $B$, so neither can be in $\olsi{B}$.
\end{feedback}
\end{prompt}
\end{problem}

\begin{problem}
Consider the statement
\[ \left(A\cap B\right)\cup\left(A\cap\olsi{B}\right) = A. \]

\begin{prompt}
\textbf{(a)} Translate the left-hand side into English, using ``and'', ``or'', and
``not''.

\begin{freeResponse}
\end{freeResponse}

\begin{feedback}
The left-hand side describes the things that are ``in $A$ and in $B$'', \emph{or}
``in $A$ and not in $B$''.
\end{feedback}
\end{prompt}

\begin{prompt}
\textbf{(b)} Is the statement true or false?

\begin{multipleChoice}
\choice[correct]{True --- the two sides are equal for any sets $A$ and $B$}
\choice{False --- there are sets $A$ and $B$ for which they differ}
\end{multipleChoice}

\begin{feedback}
It is true. Every element of $A$ either is in $B$ or is not in $B$ --- there is no
third option. So splitting $A$ into ``the part inside $B$'' and ``the part outside
$B$'' accounts for all of $A$ and nothing more, and gluing those two pieces back
together with $\cup$ returns exactly $A$.

Note also that both pieces are subsets of $A$, so the union cannot reach outside of
$A$.
\end{feedback}
\end{prompt}
\end{problem}

\begin{problem}
Suppose a particular brand of cat food has $6$ different flavors available at the
store. Euclid the cat might like all $6$ flavors, or none of them, or any other
subset of them. How many possibilities are there for the subset of flavors Euclid
likes among the $6$ available flavors?

$\answer{64}$ possible subsets.

\begin{hint}
Decide flavor by flavor: for each of the $6$ flavors, Euclid either likes it or does
not. That is $2$ choices per flavor, made independently.
\end{hint}

\begin{feedback}
A set with $n$ elements has $2^n$ subsets, since each element is independently either
in or out. Here $n=6$, so there are
\[ 2^6 = 64 \]
possible subsets of flavors --- including the empty set (Euclid likes nothing) and
the whole set (Euclid likes everything).
\end{feedback}
\end{problem}

\begin{problem}
A concert features a cellist, a flutist, a harpist, and a vocalist. Each song
features a different subset of \emph{at least two} of the featured musicians. How
many songs are in the concert if every such subset is featured once?

$\answer{11}$ songs.

\begin{hint}
Start with all $2^4$ subsets of the four musicians, then subtract the ones that are
too small to be a song: the empty set, and the four solos.
\end{hint}

\begin{feedback}
There are $2^4=16$ subsets of the four musicians in total. The subsets that do
\emph{not} have at least two members are the empty subset ($1$ of them) and the
one-musician subsets ($4$ of them), for $5$ subsets to discard. That leaves
\[ 16 - 1 - 4 = 11 \]
subsets of at least two musicians, so the concert has $11$ songs.
\end{feedback}
\end{problem}

\end{document}
